\chapter{Conclusion}
\label{conclusion}

This thesis investigated a fundamental limitation of the Ludii General Game 
Playing system: simulation-based agents involuntarily access hidden game-state 
information through unrestricted context copies, producing artificially 
inflated performance figures in imperfect-information games. This architectural 
problem renders experimental results uninterpretable and undermines the 
reliability of AI benchmarking in any Ludii game involving hidden information.

Three contributions were developed to address this problem. First, a controlled 
imperfect-information benchmark was established in the Tabletop Games Framework 
using Battleship as a case study, demonstrating that simulation consistency 
matters more than search depth in hidden-position environments. Second, the 
information leakage problem in Ludii was formally characterized and empirically 
confirmed, showing that native simulation-based agents achieve near-perfect win 
rates against non-simulating baselines purely through involuntary access to 
privileged state information. Third, a generic determinization layer was 
designed and implemented for Ludii, intercepting simulation-time context copies 
and replacing hidden information with coherent hypotheses derived from 
accumulated observations, with no modification to existing agent logic and only 
a single backward-compatible extension to the engine.

Experimental evaluation confirmed that the determinization layer achieves its 
primary objective. The re-implementation of OSLA in Ludii provided a direct 
empirical bridge between the two frameworks, with OSLA+Det achieving 87.5\% 
under constraint-based determinization, closely mirroring the 87\% observed 
in TAG under identical conditions.

\newpage

\section*{Future Work}

Several directions remain open for future work.

\paragraph{Hidden-identity games.}
The current determinization engine handles only the \texttt{who} chunk field, 
addressing hidden-position games such as Battleship. Extending it to also 
reconstruct \texttt{what} field values would enable its application to 
hidden-identity games such as Stratego, where piece locations are visible but 
ranks are concealed. This extension is conceptually straightforward but 
requires reasoning over type constraints rather than purely positional ones.

\paragraph{Richer uncertainty representations.}
The single-determinization approach adopted here inherits known limitations of 
determinization-based search. Replacing it with richer belief-state 
representations~\cite{morenville2025belief_constraints,morenville2024beliefsg} could 
improve reasoning quality under uncertainty at the cost of additional 
computational complexity. Investigating this tradeoff in the GGP setting, 
where agents must remain game-independent, constitutes a natural extension 
of this work.

\paragraph{Cross-framework comparison.}
The comparison between TAG and Ludii is currently strengthened by the 
re-implementation of OSLA+Det as a shared agent, but remains limited by the 
absence of RMHC and by differences in the internal implementations of UCT and 
MAST between the two frameworks. A more rigorous comparison would require 
either porting TAG agents to Ludii or re-implementing Ludii agents in TAG under 
identical computational conditions.

\paragraph{Porting RMHC to Ludii.}
RMHC was not ported to Ludii in this work and therefore remains absent from 
the Ludii experimental evaluation. Given its strong performance under 
constraint-based determinization in TAG, porting it would provide an 
additional data point for the cross-framework comparison and could further 
validate the generality of the determinization layer.

\paragraph{Incremental history reconstruction.}
The current history reconstruction mechanism replays the full move history 
from the Trial at each decision step, introducing an overhead that grows 
linearly with game length. Maintaining an incremental observation log updated 
after each move would substantially reduce this cost and constitutes a natural 
direction for optimization.

\paragraph{Games with complex observation structures.}
The history reconstruction currently assumes that all relevant observations 
can be inferred from the move sequence alone, which holds for Battleship but 
may not generalize to games with simultaneous moves, hidden draws, or 
probabilistic events. Extending the reconstruction mechanism to handle such 
structures would broaden the applicability of the determinization layer across 
the full Ludii game library.